\documentclass[10pt]{article}
\usepackage[margin=20mm]{geometry}
\usepackage{amsmath,amssymb,mathtools,fontspec,unicode-math,xurl,hyperref}
\IfFontExistsTF{DejaVu Serif}{\setmainfont{DejaVu Serif}}{\setmainfont{Latin Modern Roman}}
\setmathfont{latinmodern-math.otf}
\setlength{\parindent}{0pt}\setlength{\parskip}{6pt}\sloppy
\title{Initial scientific review of
Boundary integral exterior calculus}
\author{}
\date{Frozen batch M25-0248; initial review, decision pending}
\begin{document}
\maketitle
Abstract: The initial assessment is Major. The admitted two-dimensional zero-parameter case invalidates the proposed nonlocal inner products and coercivity theorem. The zero-frequency compatibility conditions are vacuous as printed and admit elementary unsolvable data. Independently, the mixed Hodge–Laplace representation has reversed boundary signs, with an explicit counterexample at a positive parameter and with smooth data on each side. Further confirmed issues concern an overstated Newton-potential range, insufficient hypotheses for the graph traces used in representation formulas, and degree-inconsistent displayed operators. Concrete counterexamples, counterevidence and bounded repairs are supplied. The positive-parameter structural idea is not refuted as a whole. Scientific decision and adjudication remain pending.\par
Article. Erick Schulz, Ralf Hiptmair and Stefan Kurz. Journal identifier: 10.4171/JEMS/1621.\par
Version examined. ETH SAM Research Report 2022-36, July 2022, whose cover specifies latest revision August 2022; 39 PDF pages. Below, p. means the printed page, and its PDF page is one greater. PDF page 1 is the cover. The publisher records online publication on 26 March 2025. The SAM directory also lists a December 2022 revision. Neither that later revision nor the journal version was acquired or compared in this review. The frozen 2025 identifier is retained, but every mathematical finding below is restricted to the actual August 2022 file.\par
\section*{Scope and mathematical mechanism}
The primary reviewer read all 39 source-text pages and actually inspected all 39 rendered source pages. Their roles, actual text and image ranges, and developer-provided model self-reports are recorded separately from the controller’s UI observation. This is an initial review, not independent acceptance.\par
The paper starts with \(D=d+\delta\) and \(-\Delta_\ell=d_{\ell-1}\delta_\ell+\delta_{\ell+1}d_\ell\). Tangential and normal traces generate two boundary de Rham complexes. Atomic potentials \(S_\lambda^\ell\) and \(D_\lambda^\ell\) are formed from the Newton potential and dual trace maps. Their commutation and jump identities lead to representations and first-kind boundary integral operators. The intended structural identification is made through bilinear forms and the nonlocal boundary inner products, rather than an unqualified equality of maps between different pivot spaces.\par
Valid potential spaces and positive boundary products are required to define the Hilbert adjoints used in the boundary Hodge structure. The distinction between \(\lambda>0\) and \(\lambda=0\), and between a maximal distributional domain and the separate \(d,\delta\) graph domains, is essential. The review checked these distinctions, endpoint degrees, signs, compatibility, and the mixed Schur complements. The full general Lipschitz trace theory, compact-manifold kernel theory and external Gaffney theorems were not independently reproved. No paper-specific supplement or code artifact was acquired in the bounded searches; this is an access/coverage limitation, not a mathematical defect.\par
\section*{Confirmed issues}
\subsection*{F1. Zero-parameter coercivity fails in the allowed dimension two (Major)}
Location. Assumption A, p. 17 (PDF 18); (3.15a) and Lemma 3.5, p. 19; Theorem 3.6, p. 20; (3.17)–(3.22), p. 21. Dependent structural conclusions occur on pp. 25–26, 30 and 34.\par
Take \(M=\mathbb R^2\), \(\Omega=B_R\), \(\ell=0\), \(\lambda=0\) and the nonzero smooth density \(h=1\) on \(\partial B_R\). These are permitted by Assumption A. For the logarithmic fundamental solution \[\Phi_0(r)=c-\frac{1}{2\pi}\log r\] with any fixed real constant \(c\), the circular logarithmic average gives \[tS_0^0h=R(2\pi c-\log R),\qquad
 (h,h)_{-1/2,0,t}=2\pi R^2(2\pi c-\log R).\] For completeness, the average of \(\log|e^{it}-e^{is}|\) is zero: it is the boundary limit of the zero mean of \(\log|1-\rho e^{i(s-t)}|\) for \(0<\rho<1\), with the integrable logarithmic singularity allowing the limit. Thus at \(R=e^{2\pi c}\) the proposed norm vanishes on a nonzero density; at larger radii it is negative. For the standard normalization \(c=0\), the unit disk already contradicts the coercive lower bound in Theorem 3.6.\par
The exterior energy explains the failure of the proof. With \(c=0\), \[S_0^0h(x)=-R\log|x| \quad (|x|>R),\qquad
 \int_{R<|x|<T}|\nabla S_0^0h|^2\,dx
   =2\pi R^2\log(T/R)\longrightarrow\infty.\] Consequently, the finite whole-space energy identity used in Lemma 3.5 is unavailable here. Boundedness of \(\Omega\) does not remove the exterior contribution at infinity. Adding a constant to the kernel does not change this energy obstruction, and a fixed constant cannot make the asserted product positive for every bounded domain.\par
There are also two distinguishable errors in the displayed coercivity proof. At \(\lambda=0\), the whole-space inhomogeneous \(H^1\) equivalence (3.21) fails: for \(V_R(x)=\chi(x/R)\), \[\|V_R\|_2^2=R^N\|\chi\|_2^2,\qquad
 \|\nabla V_R\|_2^2=R^{N-2}\|\nabla\chi\|_2^2.\] Further, the dual suprema in (3.17) and (3.22) are homogeneous of degree one in the density, whereas the displayed squared norms are homogeneous of degree two. Formula (3.18) should use a global variational identity on \(M\), not just \(\Omega\).\par
Counterevidence and repair boundary. The introduction emphasizes positive shifts, and compact manifolds are explicitly restricted to positive parameters. Neither fact cancels the Euclidean zero case admitted later. The counterexample does not refute the Yukawa case \(\lambda>0\), or establish failure of standard zero-parameter single-layer coercivity in \(N\ge3\). For \(\lambda>0\), a corrected proof starts from the unsquared dual norm, uses a bounded global trace lifting and \[a_\lambda(S_\lambda h,V)=\langle h,tV\rangle_\Gamma,\] applies Cauchy–Schwarz and the positive-parameter energy equivalence, and squares at the end. A low-dimensional zero-parameter theory instead needs explicit charge/capacity restrictions or a changed normalization and space, with corresponding new proofs. Since the boundary Hilbert products underlie the main adjoint and Fredholm arguments, this is a substantive scope defect.\par
\subsection*{F2. The zero-frequency compatibility kernels are interchanged (Major)}
Location. (2.4a,b) and (2.5a,b), p. 8 (PDF 9); the compatibility discussion and (2.9b), p. 9.\par
By definition, \(nV=0\) on \(\ker D_n\) and \(tV=0\) on \(\ker D_t\). Hence the printed conditions \[\langle g,nV\rangle_\Gamma=0\quad(V\in\ker D_n),\qquad
 \langle h,tV\rangle_\Gamma=0\quad(V\in\ker D_t)\] are both automatic.\par
Let \(\Omega\) be the unit disk in \(\mathbb R^2\) and \(\kappa=0\). Prescribe only the tangential degree-one component \(g_1=d\theta\), with all other components zero. Any \(DU=0\) in the declared graph domain must satisfy \(dU_1=0\). Weak Stokes gives \[0=\int_\Omega dU_1=\int_{\partial\Omega}tU_1,
 \qquad
 \int_{\partial\Omega}g_1=2\pi,\] so no solution exists, although (2.5a) is satisfied. Likewise, prescribing only the normal scalar component \(h_0=1\) is incompatible with \(\delta U_1=0\), because its total outward flux is \(2\pi\) instead of zero.\par
The correct tangential compatibility tests against \(\ker D_t\), and the correct normal compatibility tests against \(\ker D_n\). The constant volume form in \(\ker D_t\) detects the first obstruction; the constant scalar in \(\ker D_n\) detects the second. The right-hand side of (2.9b) should also involve \(tV\), as in (2.8b), rather than the printed \(nV\).\par
Impact and limits. These examples disprove the asserted solvability of the original zero-frequency problems, including on an orthogonal complement. A solvable augmented system with a nonzero harmonic multiplier does not solve the original homogeneous equation. The kernel-label correction is short, but it changes the actual compatibility criterion and is central to the stated BVP scope. The nonzero-\(\kappa\) problems are not counterexamples.\par
\subsection*{F3. The mixed representation has reversed boundary signs (Major)}
Location. Proposition 4.6 and its proof, p. 31 (PDF 32); potential definitions, p. 32. Compare the general scheme on p. 22 and the jump convention on p. 18.\par
Take \(M=\mathbb R^3\), \(\Omega=B_1\), \(\ell=1\), \(\lambda=1\), and \[U=(\mathbf1_{B_1},0),\qquad F=(-\mathbf1_{B_1},0).\] This pair is smooth on each side, belongs to the required piecewise graph domains, and satisfies the printed piecewise mixed equation. Put \[q=N_1^0\mathbf1_{B_1},\qquad P=D_1^1 1,\] where the subscript here denotes the parameter and the superscript the degree. Since \(d\mathbf1_{B_1}=-n_1'1\) distributionally, \[P=-dq,\qquad \delta P=-\mathbf1_{B_1}+q.\] With the paper’s exterior-minus-interior jump, \([tU_0]=-1\) and all normal data vanish. The printed proposition therefore applies \[A=\begin{pmatrix}-1&\delta\\d&1\end{pmatrix}
 \quad\hbox{to}\quad
 \binom{-q}{+P}.\] Its first component is \(2q-\mathbf1_{B_1}\), instead of \(\mathbf1_{B_1}\).\par
This is an explicit pointwise discrepancy away from the boundary. The radial Yukawa potential is \[q(r)=
 \begin{cases}
 1-\dfrac2e\,\dfrac{\sinh r}{r},&0<r<1,\\[4pt]
 \dfrac{e^{-(r+1)}}r,&r>1.
 \end{cases}\] With \(q(0)=1-2/e\), it is regular at zero, is \(C^1\) at \(r=1\), solves \((-\Delta+1)q=1\) inside and \(0\) outside, and decays at infinity. At \(r=2\), the printed first component is \(e^{-3}>0\), while \(U_0=0\). These identities were also checked symbolically from the displayed radial functions.\par
Correct form and dependency. With the potential definitions retained on p. 32, the representation must read \[U=N[M]F-\operatorname{SL}[M]\,[T_M^nU]+\operatorname{DL}[M]\,[T_M^tU].\] In the example, the corrected formula applies \(A\) to \((-q,-P)\) and recovers \(\mathbf1_{B_1}\). In general, distributional differentiation gives the normal source vector \(-(t'_{\ell-1}[n_\ell U_\ell],t'_\ell[n_{\ell+1}d_\ell U_\ell])\) and the tangential source vector \((0,n'_\ell[t_{\ell-1}U_{\ell-1}]+\delta_{\ell+1}n'_{\ell+1}[t_\ell U_\ell])\). Applying the stated inverse gives exactly these corrected signs.\par
A related display on p. 28 (PDF 29) gives the normal jump of \(\operatorname{SL}[\Delta]\) as the positive input vector. It must be its negative, consistent with (3.14b). For the scalar Laplace single layer of constant density on the unit sphere, the potential is \(1\) inside and \(1/r\) outside, so the normal-derivative jump is \(-1\).\par
Counterevidence and repair boundary. The outer differential factor in Proposition 4.6 is present: actual image inspection shows that the large parenthesis encloses every term. An earlier missing-factor candidate was rejected. The confirmed sign error is different and survives positive parameters and regular graph traces. Reverse the two boundary signs in the proposition and its proof and recheck the Calderón equations with the fixed convention. The structural weak forms of the separately defined BIOs are not thereby disproved. Nevertheless, a central representation theorem used to connect the BVP and BIE is false as printed, which warrants Major classification.\par
\subsection*{F4. The Newton-potential range is too strong across its source (Minor)}
Location. The graph-space definitions on p. 3 (PDF 4), (3.3) on p. 17, and Lemma 3.3 on p. 18.\par
The asserted mapping from \(H^{-1}_{\mathrm{comp}}\) includes \(H_{\mathrm{loc}}\Lambda(\Delta,M)\), whose definition requires the relevant second derivatives to be locally square-integrable. In \(M=\mathbb R\), with \(\lambda=1\), \[f=\delta_0\in H^{-1}_{\mathrm{comp}}(\mathbb R),\qquad
 N_1^0f(x)=\tfrac12e^{-|x|}.\] Membership of \(\delta_0\) in \(H^{-1}\) follows from the finite integral \(\int_\mathbb R(1+\xi^2)^{-1}\,d\xi\). The potential is in \(H^1\), but its first derivative has jump \(-1\), and \[(N_1^0f)''=N_1^0f-\delta_0\notin L^2_{\mathrm{loc}}.\] In higher dimensions a nonzero smooth hypersurface density \(t'h\in H^{-1}\setminus L^2\) gives the same contradiction by applying \(-\Delta+\lambda\) to the proposed graph-regular potential.\par
Impact and repair. Keep the global \(H^1_{\mathrm{loc}}\) conclusion and restrict second-order regularity to the complement of the source support, or strengthen the source class. For the intended boundary graph densities, \(\delta S h=S\delta h\) gives sidewise \(d\delta S h\in L^2\); the homogeneous side equation then gives \(\delta dS h\in L^2\). Thus the surplus across-source regularity can be removed without changing the intended sidewise potential calculus. This explicit closure, rather than mere presumed repairability, supports the Minor classification.\par
\subsection*{F5. The representation hypotheses do not provide the declared traces (Minor)}
Location. Proposition 4.3 and its proof, pp. 23–24 (PDF 24–25); Proposition 4.5, p. 27; graph domains on pp. 3–4.\par
In the unit disk take the holomorphic series \[f(z)=\sum_{k\ge1}k^{-1/2}z^k=u+iw\] and extend its real and imaginary parts by zero outside. Orthogonality gives \[\int_{B_1}|u|^2=\frac{\pi}{2}\sum_{k\ge1}\frac1{k(k+1)}
 =\frac{\pi}{2},
 \qquad
 \int_{B_1}|\nabla u|^2=\pi\sum_{k\ge1}1=\infty.\] The scalar function is piecewise harmonic and satisfies the printed assumptions of Proposition 4.5 with \(F=\lambda u\), even for \(\lambda>0\). But its boundary Fourier series is not in \(H^{1/2}(S^1)=H^{-1/2}\Lambda^0(d,S^1)\), the declared tangential graph-trace space.\par
For Proposition 4.3 set \(U=u-w\,dx\wedge dy\). The Cauchy–Riemann equations give \[\delta(-w\,dx\wedge dy)=-du,\qquad DU=0\] on each side. For any \(\kappa\ne0\), take \(F=i\kappa U\). The compactly supported \(L^2\) and piecewise equation assumptions hold, but \(dU\) and \(\delta U\) are not separately \(L^2\). Hence \[U\notin H\Lambda(D)
 =H\Lambda(d)\cap H\Lambda(\delta),\] contrary to the regularity implication used to justify the independent traces.\par
Impact and repair boundary. Require the appropriate piecewise componentwise graph domains in these propositions. Those are precisely the domains of the intended BVPs and of the Green identities used in the proofs. Alternatively, explicitly develop maximal-domain weak traces and their larger distribution spaces. This finding does not deny the existence of such weak traces or disprove a suitably extended representation theorem. Restricting to the already-declared BVP domains closes this domain issue, while the regular example in F3 remains a counterexample to its printed signs.\par
\subsection*{F6. Degree labels in several operator displays are inconsistent (Minor)}
Location. (1.5), p. 3; the tangential product below (2.16), p. 12; (2.21), p. 14; (2.25), p. 15; (2.28a), p. 16; (3.7), p. 17; (3.15b), p. 19; (4.16b), (4.19b) and Fig. 4, pp. 30, 34, 36.\par
All corrections follow from the stated maps \(d_j:\Lambda^j\to\Lambda^{j+1}\), \(\delta_j:\Lambda^j\to\Lambda^{j-1}\), \(\star_j:\Lambda^j\to\Lambda^{N-j}\).\par
- The full-star matrix must satisfy \((\star U)_i=\star_{N-i}U_{N-i}\). Its top-right block is \(\star_N\), and its bottom-left block is \(\star_0\), the reverse of (1.5). Already \(N=1\) detects the domain mismatch.\par
- In (3.7), \(D_\lambda^\ell=N_\lambda^\ell n_\ell'\), not \(N_\lambda^\ell n_{\ell-1}'\). In (3.15b), the normal pairing uses \(n_{\ell+1}D_\lambda^{\ell+1}\), not \(n_{\ell+1}D_\lambda^\ell\).\par
- The upper-right block of (4.16b), (4.19b) and Fig. 4 is \(-d_{\ell-1}^*\); the analogous bulk block in (2.28a) is \(d_{\ell-1}^*\). The printed \(d_\ell^*\) has the wrong input and output degrees.\par
- In (2.21), \(d_{\ell+1}\delta_\ell\) must be \(d_{\ell-1}\delta_\ell\). In (2.25) and Fig. 4, \(d_\ell V_{\ell-1}\) must be \(d_{\ell-1}V_{\ell-1}\). The tangential product below (2.16) must be ordered \((\ell-1,\ell)\), consistently with (1.16).\par
The adjoint error can be tested against the preceding weak form. On \(S^1\), choose \(\ell=1\), \(\lambda>0\), \(g_1=d_0(\sin\theta)\), \(v_0=\sin\theta\), and \(g_0=v_1=0\). Formula (4.15) gives \(-\|g_1\|_{-1/2,\lambda,n}^2<0\). The printed block \(d_1^*\) is not defined on a one-form; extending that zero top-degree operator by zero cannot produce this response. The corrected matrix is \[W=
 \begin{pmatrix}
 I&-d_{\ell-1}^*\\
 -d_{\ell-1}&-d_\ell^*d_\ell-\lambda I
 \end{pmatrix},\] whose Schur complement is \(-(d_{\ell-1}d_{\ell-1}^*+d_\ell^*d_\ell+\lambda I)\). The intended positive-parameter Hodge–Laplace structure is thereby restored. The correctly typed weak forms are counterevidence against interpreting these display errors as a failure of every underlying structural argument.\par
\section*{Positive checks, limitations and recommendation}
Actual source images show the necessary outer grouping in both (4.8) and Proposition 4.6, and the relevant conjugations and absolute values in the complex coercivity argument. These text-extraction candidates were explicitly excluded. In a valid nonlocal Hilbert complex, self-adjointness gives \[\|(T-i\kappa)x\|^2=\|Tx\|^2+\kappa^2\|x\|^2,\] supporting the nonzero imaginary-shift argument. The corrected mixed matrix has the expected positive Laplace Schur complement. In a legitimate zero-parameter complex, the kernels of the two mixed operators correspond to harmonic degrees \(\ell-1\) and \(\ell\), respectively; the review does not reject the abstract Hilbert-complex theory.\par
The required substantive revisions to this frozen version are an explicit valid parameter/domain range for the nonlocal products, the correct zero-mode compatibility conditions, and a consistent representation/jump-sign derivation. The graph-domain and degree corrections should then be propagated through the displayed statements and summary figures. A corrected positive-parameter theory is plausible and is supported by several checked identities, but that observation cannot certify the printed statements or the unexamined journal version.\par
The evidence package contains per-finding hypotheses, exact source pages, proofs, counterevidence and repair limits; a reviewer-authored symbolic check; actual image-inspection records; the current TeX/PDF/compiler/dependency hashes; and retained earlier candidates and tool failures. No author code was executed. Final scientific decision and human review remain pending.\par
9 E. Schulz, R. Hiptmair and S. Kurz, Boundary Integral Exterior Calculus, ETH SAM Research Report 2022-36, revision August 2022, frozen 39-page source. EMS Press, article metadata for Boundary integral exterior calculus, https://ems.press/journals/jems/articles/14298685, accessed 27 October 2026. ETH SAM, report 2022-36 and revision directory, https://math.ethz.ch/sam/research/reports/.ethz\_search.html?id=1024, accessed 27 October 2026. Metadata only for later versions.\par
\end{document}
